\clearpage
\section{sysfs：用 cat 看硬體狀態}
\label{sec:sysfs}

\begin{oneline}
六個 cycle counter、STATUS 和 NBLK 各對應 sysfs 裡的一個唯讀檔案，\code{cat} 一下就能看到，不必寫程式，也不必開裝置。
\end{oneline}

\begin{before}
counter 只能由 app 在跑完一幀後讀出來（第 283–288 行，經由 \code{mmap} 進來的 \code{/dev/uio4}）：

\snip{uio_counters}

沒有 app 在跑時，想看一眼硬體狀態，就得另外寫程式去 \code{mmap} \code{/dev/uio4}，或用 \code{devmem} 直接讀實體位址。
\end{before}

\begin{tech}
\begin{itemize}
\item \code{DEVICE_ATTR_RO(name)}：宣告一個唯讀的 sysfs 屬性，kernel 讀檔時呼叫 \code{name_show()}。
\item \code{sysfs_emit()}：把一個值格式化寫進 sysfs 的輸出緩衝區，保證不會寫超過一頁（4096 bytes）。
\item \code{ATTRIBUTE_GROUPS()}、\code{driver.dev_groups}：把一組屬性交給 driver core，綁定成功後自動建立檔案，解除綁定前自動移除。
\item \texttt{\#\#}：C 前處理器的「把兩個記號接起來」，用一個巨集產生八個函式。
\end{itemize}
\end{tech}

\begin{gain}
\begin{itemize}
\item driver 綁定期間，任何時候都能 \code{cat} 硬體狀態，不用開裝置（板上輸出見下）。
\item 一個檔案一個值，這是 sysfs 的慣例，腳本好解析。
\item 檔案的建立與移除交給 driver core，和 \code{probe()} / \code{remove()} 不會有先後的競爭問題。
\end{itemize}
\end{gain}

\subsection{實際例子：程式碼與板上輸出}

\textbf{一個巨集產生一個屬性。}\code{TWN_REG_ATTR(name, off, fmt)} 展開成一個 \code{name_show()} 函式（讀一個暫存器、照格式印出），再加上 \code{DEVICE_ATTR_RO(name)}：

\snip{drv_reg_attr}

用它宣告八個屬性：

\snip{drv_reg_attr_use}

再列成一張表，\code{ATTRIBUTE_GROUPS(twn)} 產生 \code{twn_groups}，交給第 1 節 \code{twn_driver} 的 \code{.dev_groups}：

\snip{drv_attrs}

\textbf{板上：跑完一幀後讀出來。}

\boardfile{A3_driver_state.txt}{16-25}

怎麼讀這些數字（PL 時脈 240 MHz）：
\begin{itemize}
\item \code{cycles_total} 是 17 個 block 執行期間的 cycle 總和（\code{rtl/fuse/twn_blockseq.v} 只在 block 執行與收尾時計數，block 之間的空檔不算）。\cyclesTotal{} cycles $\div$ 240 MHz = 4.13 ms，是 START 到 done 的下限；實際的 START 到 done 時間看 \code{TWN_IOC_BB_RUN} 回傳的 \code{ns}。
\item \code{cycles_array} 976,076 cycles，占 98.5\%：三值運算陣列幾乎一直在工作。
\item \code{cycles_stall} 只有 167 cycles：記憶體幾乎沒有讓運算停下來。
\item \code{status} = \code{0x00010043}，欄位見下表。
\end{itemize}

\begin{center}
\small
\begin{tabularx}{\textwidth}{l l X}
\toprule
位元 & 名稱 & 這次的值與意思 \\
\midrule
0 & \code{d_valid} & 1：描述子表通過檢查 \\
1 & \code{tab_checked} & 1：檢查已經跑過 \\
6:2 & \code{cur_block} & 16：最後執行的 block 編號（0 到 16） \\
10:7 & \code{err_code} & 0：表格沒有錯；第 6 節被拒絕的表，這裡是 1（\code{SLOT_RANGE}） \\
15:11 & \code{err_block} & 0：出錯的 block 編號 \\
16 & \code{ovr_sticky} & 1：見下段 \\
\bottomrule
\end{tabularx}
\end{center}

\textbf{STATUS bit 16：一個 driver 之外的發現。}bit 16 是 ring overrun 的黏著旗標：PL 裡某個列緩衝區在滿的時候被寫入，或被讀的時候已經不是要的那一列，就會設成 1（\code{rtl/fuse/twn_fuse_top.v} 第 755–770 行；只記錄、不停機），只有 PL reset 會清掉。板上量到：

\boardfile{A1_fresh_boot_uio_dtb.txt}{9-10}
\boardfile{C1_rollback_to_uio_dtb.txt}{10-21}

第一個框是剛開機、什麼都還沒跑：0。第二個框是回退到原本的 dtb、重開機後，driver 完全沒有載入（\code{lsmod} 數到 0），只用 v1.0.0 的 UIO 版本跑 4 張 fixture：跑之前是 0，跑完變成 \code{0x00010043}。所以這個旗標是 bitstream 本身在正常執行時就會設起來的，和 driver 無關；輸出照樣逐位元正確。原因還沒查明，記錄在 \code{results/board_results.json} 這次那一筆的 \code{note} 欄位，留待之後追查。

\code{ls} 這個目錄，可以看到 driver 加的八個檔案，以及 driver core 本來就有的檔案（\code{driver}、\code{driver_override}、\code{of_node}、\code{modalias} 等）：

\boardfile{A3_driver_state.txt}{26-43}

\subsection{原理}

\textbf{sysfs 是什麼。}sysfs 是 kernel 把內部物件（裝置、driver、bus）攤開成檔案樹的檔案系統，掛在 \code{/sys}。每個 \code{struct device} 都是一個目錄。讀其中一個屬性檔時，kernel 呼叫對應的 \code{show()}，把它寫進緩衝區的字串當成檔案內容。

\textbf{巨集展開後長什麼樣。}以第 \lineof{drv_attr_nblk} 行的 \code{TWN_REG_ATTR(nblk, R_NBLK, ...)} 為例，\texttt{name\#\#\_show} 會被接成 \code{nblk_show}：

\begin{lstlisting}[numbers=none,xleftmargin=0pt,title={\scriptsize 巨集展開後的樣子（示意，不是原始檔的內容）}]
static ssize_t nblk_show(struct device *dev, struct device_attribute *a, char *buf)
{
	struct twn_dev *td = dev_get_drvdata(dev);

	return sysfs_emit(buf, "%u" "\n", rd(td, R_NBLK));
}
static DEVICE_ATTR_RO(nblk);   /* defines dev_attr_nblk, mode 0444, show = nblk_show */
\end{lstlisting}

\code{DEVICE_ATTR_RO(nblk)} 依命名規則去找 \code{nblk_show}，並產生 \code{dev_attr_nblk} 這個變數；\code{twn_attrs[]} 裡放的正是 \code{&dev_attr_nblk.attr}。

\textbf{為什麼用 \code{.dev_groups}，不在 \code{probe()} 裡自己建檔。}Linux 6.1 的 \code{drivers/base/dd.c}：\code{really_probe()} 在第 658 行呼叫 \code{device_add_groups()}，時間點是 \code{probe()} 成功之後；第 545 行則在呼叫 \code{remove()} 之前先移除。所以 \code{show()} 被呼叫時，\code{dev_get_drvdata()} 一定拿得到 \code{probe()} 設好的資料，不會有「檔案已經出現、資料還沒準備好」的空窗。

\textbf{什麼該放 sysfs，什麼該放 ioctl。}
\begin{itemize}
\item sysfs：一個檔案一個值、給人看或給腳本讀的狀態與設定。
\item ioctl：有結構的參數、一次做完一整段操作、需要回傳錯誤碼的動作。
\item debugfs（這次沒用）：只給開發者除錯用、格式不保證穩定的資訊。
\end{itemize}

\subsection{面試會問}

\begin{qa}{除了 ioctl，你怎麼讓使用者看到硬體狀態？}
我用 sysfs。六個 cycle counter、STATUS 和 NBLK 各是一個唯讀檔案，用一個巨集產生 \code{show} 函式加 \code{DEVICE_ATTR_RO}，再用 \code{ATTRIBUTE_GROUPS} 交給 platform driver 的 \code{.dev_groups}。driver core 會在 probe 成功後建檔、remove 前移除，所以不會讀到還沒初始化的資料。板上跑完一幀，\code{cat cycles_total} 是 \cyclesTotal{} cycles，這是 17 個 block 執行期間的總和，約 4.13 ms，運算陣列忙碌 98.5\%。原則是：sysfs 一個檔一個值、給人和腳本看；有結構、要回錯誤碼的動作用 ioctl。
\end{qa}

\subsection{關鍵函式速查}

\begin{center}
\small
\begin{tabularx}{\textwidth}{l l X}
\toprule
名稱 & 位置 & 做什麼 \\
\midrule
\code{TWN_REG_ATTR} & \code{twn_fuse.c}:\lineof{drv_reg_attr} & 一個巨集產生 \code{show()} 與 \code{DEVICE_ATTR_RO} \\
八個屬性 & \code{twn_fuse.c}:\lineof{drv_reg_attr_use} & 六個 counter、\code{status}、\code{nblk} \\
\code{twn_attrs[]} & \code{twn_fuse.c}:\lineof{drv_attrs} & 屬性表與 \code{ATTRIBUTE_GROUPS(twn)} \\
\code{really_probe()} 呼叫 \code{device_add_groups} & kernel \code{drivers/base/dd.c}:658 & probe 成功後建立 \code{dev_groups} \\
\bottomrule
\end{tabularx}
\end{center}
